% !TeX program = lualatex
% halo:begin
% {
%   "version": 1,
%   "publish": true,
%   "course": "higher-algebra-i",
%   "section": "1.3",
%   "title": "高等代数（一）1.3：对换与行列式性质",
%   "categories": ["study-notes"],
%   "tags": ["mathematics"],
%   "excerpt": "对换与排列奇偶性、行列式性质，以及第二章习题4、6、8。"
% }
% halo:end
% 课堂日期：2026-09-22。对照王萼芳、石生明修订《高等代数》第五版第二章。
% 按一行（列）展开迁至9月24日笔记；保留本文件路径以维持已有发布链接。
% 2026-10-05补证：核对北大《高等代数》第六版第二章，印刷页35--43（PDF页40--48）。
% 保留前篇按列选取元素的定义，并补证它与教材按行写法的等价性。
\RequirePackage{iftex}
\RequireLuaTeX
\documentclass[UTF8, fontset=fandol, a4paper, 12pt]{ctexart}

\usepackage[margin=2.5cm]{geometry}
\usepackage{amsmath, amssymb, amsthm}
\usepackage{enumitem}
\usepackage{fancyhdr}
\usepackage[hidelinks]{hyperref}

\theoremstyle{definition}
\newtheorem{theorem}{定理}[section]
\newtheorem{proposition}[theorem]{命题}
\renewcommand{\proofname}{证明}
\setlist[enumerate]{leftmargin=2em, itemsep=0.2em, topsep=0.3em}
\setlength{\headheight}{16pt}
\pagestyle{fancy}
\fancyhf{}
\fancyhead[L]{\small 高等代数（一）}
\fancyhead[R]{\small 1.3\quad 对换与行列式性质}
\fancyfoot[C]{\thepage}
\renewcommand{\headrulewidth}{0.3pt}
\raggedbottom

\title{对换与行列式性质}
\author{Yuchen Zhang}
\date{课堂日期：2026年9月22日}
\makeatletter
\renewcommand{\maketitle}{%
  \begin{center}
    {\LARGE\bfseries \@title\par}\vspace{0.6em}
    \@author\par\vspace{0.2em}
    {\small \@date\par}
  \end{center}
}
\makeatother

\begin{document}
\maketitle

行列式的排列展开式为每一种元素选择规定了正负号。
要理解交换行、列为什么使行列式变号，我们先从排列中两个数的对换看起。

\section{对换与排列奇偶性}

\begin{theorem}
排列中任意两个数对换，排列的奇偶性改变。
\end{theorem}
\begin{proof}
先看相邻两数的对换。其余每个数都在它们的同一侧，
所以它与这两个数形成的逆序总数不变；只有被交换两数之间的逆序状态反转。
因此逆序数增加或减少 $1$，奇偶性改变。

若两数之间隔着 $s$ 个数，先把左边的数向右移动 $s+1$ 个位置，
再把原来右边的数向左移动 $s$ 个位置。
这样两个数交换了位置，中间各数恢复原来的次序，共作了 $2s+1$ 次相邻对换。
每次都改变奇偶性，而 $2s+1$ 为奇数，故原对换也改变奇偶性。
\end{proof}

\begin{proposition}
在任意一个将排列变为自然顺序的对换过程中，对换次数与原排列的逆序数同奇偶。
\end{proposition}
\begin{proof}
只要排列还不是自然顺序，就存在相邻且左边大于右边的两个数。
交换这两个数，逆序数减少 $1$。逆序数是非负整数，故重复此操作必在有限步后得到自然顺序。

对任意还原过程，设共作了 $m$ 次对换。每次对换都改变奇偶性，
而自然顺序的逆序数为 $0$，故 $m$ 与原排列的逆序数同奇偶。
\end{proof}

按列选出的元素也可以按行重排，两种写法中的指标排列互为逆排列。
因此，在改换写法之前，先比较原排列与逆排列的逆序数。

\begin{proposition}[逆排列的逆序数]
对任意排列 $\sigma\in S_n$，有
\[
 \tau(\sigma^{-1})=\tau(\sigma).
\]
\end{proposition}
\begin{proof}
任取 $\sigma$ 的一个逆序对 $(i,j)$，即 $i<j$ 且 $\sigma(i)>\sigma(j)$。
在逆排列中，位置 $\sigma(j)<\sigma(i)$，而这两个位置上的数满足
\[
 \sigma^{-1}(\sigma(j))=j>i=\sigma^{-1}(\sigma(i)).
\]
所以 $(\sigma(j),\sigma(i))$ 是 $\sigma^{-1}$ 的一个逆序对。
反过来，对逆排列的逆序对应用同样的对应，就恢复原来的 $(i,j)$。
这给出了两种排列的逆序对之间的一一对应，故逆序总数相等。
\end{proof}

\section{行列式的性质}
以下设 $A=(a_{ij})$ 为数域 $P$ 上的 $n$ 阶方阵。

\subsection{按行与按列的排列展开}

上一篇按列的顺序定义行列式，每项的行指标构成一个排列。
借助逆排列，我们也可以把同一项的因子按行重排：
\[
 \prod_{j=1}^n a_{\sigma(j),j}
 =\prod_{i=1}^n a_{i,\sigma^{-1}(i)},\qquad i=\sigma(j).
\]
令 $\pi=\sigma^{-1}$。由逆序数相等，且 $\pi$ 也遍历全部排列，得
\[
 \det A
 =\sum_{\sigma\in S_n}(-1)^{\tau(\sigma)}\prod_{j=1}^n a_{\sigma(j),j}
 =\sum_{\pi\in S_n}(-1)^{\tau(\pi)}\prod_{i=1}^n a_{i,\pi(i)}.
\]
两式选取的是同一组元素，区别只在于按列还是按行写出它们。

\subsection{转置与三角形行列式}

\begin{proposition}
对任意方阵 $A$，有 $\det A^{\mathsf T}=\det A$。
若 $A$ 为上三角形或下三角形矩阵，则
\[
  \det A=\prod_{i=1}^{n}a_{ii}.
\]
\end{proposition}
\begin{proof}
对 $A^{\mathsf T}$ 使用按列展开式，得到
\[
 \det A^{\mathsf T}
 =\sum_{\sigma\in S_n}(-1)^{\tau(\sigma)}\prod_{j=1}^n a_{j,\sigma(j)}.
\]
右端正是 $A$ 的按行展开式，故 $\det A^{\mathsf T}=\det A$。

若 $A$ 为上三角形矩阵，则 $i>j$ 时 $a_{ij}=0$。
因此按行展开式中的非零项必须满足 $\pi(i)\ge i$（$i=1,\ldots,n$）。
但排列的指标总和不变，即 $\sum_i\pi(i)=\sum_i i$，
故这些不等式只能全部取等号，否则左边总和会严格大于右边。
所以只有自然排列的一项可能非零，且其符号为正，得到 $\det A=\prod_i a_{ii}$。
下三角形情形由转置得到。
\end{proof}

\subsection{多重线性与交替性}

\begin{theorem}
行列式分别对每一行、每一列线性。交换两行（列），行列式变号。
\end{theorem}
\begin{proof}
固定其余各行，把第 $r$ 行写成 $s\boldsymbol u+t\boldsymbol v$，其中 $s,t\in P$，
$\boldsymbol u,\boldsymbol v$ 为 $n$ 维行向量，所得矩阵记为 $B$。
用 $A_{\boldsymbol u}$、$A_{\boldsymbol v}$ 分别表示将 $A$ 的第 $r$ 行换为
$\boldsymbol u$、$\boldsymbol v$ 后的矩阵。按行展开并使用分配律，得
\begin{align*}
 \det B
 &=\sum_{\pi\in S_n}(-1)^{\tau(\pi)}
   (su_{\pi(r)}+tv_{\pi(r)})\prod_{i\ne r}a_{i,\pi(i)}\\
 &=s\det A_{\boldsymbol u}+t\det A_{\boldsymbol v}.
\end{align*}
这就证明了对第 $r$ 行的线性；关键在于每项只含这一行的一个元素。

再设 $B$ 由 $A$ 交换第 $r,s$ 行得到。为了跟踪行指标，使用按列展开式。
对任意排列 $\sigma$，把其中的 $r,s$ 两个值互换，所得排列记为 $\sigma'$，则
\[
 \prod_{j=1}^n b_{\sigma(j),j}=\prod_{j=1}^n a_{\sigma'(j),j},
 \qquad (-1)^{\tau(\sigma')}=-(-1)^{\tau(\sigma)}.
\]
这里两个乘积选取同一组元素，而对换使符号相反。
由于 $\sigma\mapsto\sigma'$ 是一一对应，对全部排列求和就得到 $\det B=-\det A$。
列的结论由转置性质得到。
\end{proof}

\begin{proposition}
由多重线性与交换行（列）变号，有以下结论：
\begin{enumerate}
 \item 一行（列）乘 $c$，行列式乘 $c$；有零行（列）时，行列式为 $0$；
 \item 两行（列）相同或成比例，行列式为 $0$；
 \item 把一行（列）的 $c$ 倍加到另一行（列），行列式不变；
 \item 某一列可由其余列线性表示时，行列式为 $0$，行的情形同理。
\end{enumerate}
\end{proposition}
\begin{proof}
第一条是线性中的数乘性质，取 $c=0$ 即得零行的结论。
若两行相同，交换后矩阵不变，而行列式变号，故 $D=-D$。
由于 $P\subseteq\mathbb C$，有 $2\ne0$，从而 $D=0$。
若一行是另一行的倍数，提出该倍数后便归结为两行相同的情形。

这里用 $D(\ldots,\boldsymbol u,\ldots,\boldsymbol v,\ldots)$ 表示以括号内的向量依次作为各行的行列式，则
\begin{align*}
 &D(\ldots,\boldsymbol u,\ldots,\boldsymbol v+c\boldsymbol u,\ldots)\\
 &\quad=D(\ldots,\boldsymbol u,\ldots,\boldsymbol v,\ldots)
       +cD(\ldots,\boldsymbol u,\ldots,\boldsymbol u,\ldots)\\
 &\quad=D(\ldots,\boldsymbol u,\ldots,\boldsymbol v,\ldots).
\end{align*}
第二项中的行列式有两行相同，因而为零，这就证明了倍加不变。

最后，若第 $j$ 列满足 $\boldsymbol\alpha_j=\sum_{k\ne j}c_k\boldsymbol\alpha_k$，
则沿第 $j$ 列使用线性，得到若干个含两列相同的行列式，故每项都为零。
行的结论可同样证明。
\end{proof}

\subsection{这些性质怎样决定排列展开式}

以下重新按列向量写行列式，记 $A$ 的第 $j$ 列为 $\boldsymbol\alpha_j$。
将每一列按标准单位列向量拆开，有 $\boldsymbol\alpha_j=\sum_{i=1}^n a_{ij}\boldsymbol e_i$。
逐列使用线性，得
\[
 \det A=\sum_{i_1=1}^n\cdots\sum_{i_n=1}^n
 a_{i_1 1}\cdots a_{i_n n}\det(\boldsymbol e_{i_1},\ldots,\boldsymbol e_{i_n}).
\]
如果两个行指标相同，对应的两个单位列向量也相同，因而这一项为零。
剩下的指标恰好构成排列 $\sigma(1),\ldots,\sigma(n)$。
将这些列交换到自然顺序，对换次数与 $\tau(\sigma)$ 同奇偶；
又由三角形行列式的公式，有 $\det(\boldsymbol e_1,\ldots,\boldsymbol e_n)=1$，因此
\[
 \det(\boldsymbol e_{\sigma(1)},\ldots,\boldsymbol e_{\sigma(n)})
 =(-1)^{\tau(\sigma)}.
\]
代回上式，就得到原来的排列定义。
所以多重线性、交换两列变号与标准单位向量组的取值 $1$，共同决定了行列式的排列展开式。
任何满足这三个条件的运算，都能按相同的过程展开，得到相同的值。

\section{课后作业}

\subsection*{第二章习题 4}
判定排列 $n(n-1)\cdots21$ 的逆序数，并讨论它的奇偶性。
\begin{proof}[解]
任意两数均构成逆序，故逆序数为
\[
  (n-1)+(n-2)+\cdots+1=\frac{n(n-1)}2.
\]
当 $n\equiv0,1\pmod4$ 时为偶排列；
当 $n\equiv2,3\pmod4$ 时为奇排列。
\end{proof}

\subsection*{第二章习题 6}
求六阶行列式展开式中
$a_{23}a_{31}a_{42}a_{56}a_{14}a_{65}$ 与
$a_{32}a_{43}a_{14}a_{51}a_{66}a_{25}$ 两项的符号。
\begin{proof}[解]
按行指标排列因子，第一项的列指标排列为 $431265$，
逆序数为 $3+2+0+0+1+0=6$；第二项的列指标排列为 $452316$，
逆序数为 $3+3+1+1+0+0=8$。两项均带正号。
\end{proof}

\subsection*{第二章习题 8（1）（3）}
按定义计算两个 $n$ 阶行列式：第（1）小题的非零元素为
$a_{i,n+1-i}=i$（$1\le i\le n$）；第（3）小题的非零元素为
$a_{i,n-i}=i$（$1\le i<n$）及 $a_{nn}=n$。
\begin{proof}[解]
第（1）小题仅有列指标排列 $n(n-1)\cdots21$ 的一项非零，
其逆序数为 $\binom n2$，故
\[
  D_1=(-1)^{\binom n2}n!.
\]
第（3）小题仅有列指标排列 $(n-1)(n-2)\cdots1n$ 的一项非零，
其逆序数为 $\binom{n-1}{2}$，故
\[
  D_3=(-1)^{\binom{n-1}{2}}n!.
\]
\end{proof}

\end{document}
